% 英文源分章 PDF 第 8—11 页；第 9.5—9.6 节。
\section{CV-TMLE}

$\Psi_{P_{n,v^c}}(P_0)$ 的特定于 $v$ 的 TMLE 仅基于大小为 $n/V$ 的样本 $P_{n,v}$。幸运的是，可以对上一节的程序作出修改，在 Hubbard 等（2016）定理 1 所采用的同样统计假设下，同时实现两个表面上相互冲突的目标：(1) 在构造数据自适应参数估计量时，利用更多数据来估计其中用到的数据生成分布；(2) 不因过拟合而增加偏差。CV-TMLE 的理论及一般框架（van der Laan 和 Rose，2011，第 27 章）表明，例如，CV-TMLE 比标准 TMLE 更稳健：它能够在更大的统计模型中保证渐近抽样分布的有关结论。此外，CV-TMLE 也可用于估计本章着重讨论的这类数据自适应参数。例如，它尤其适合于一方面用数据估计最优处理规则（Luedtke 和 van der Laan，2016b），另一方面估计采用该规则对结局均值的影响（van der Laan 和 Luedtke，2014；Luedtke 和 van der Laan，2016a）。在讨论变量重要性应用中这一特定的数据自适应参数的估计之前，我们先按照 van der Laan 和 Luedtke（2014），对适用于一般数据自适应参数的 CV-TMLE 作出一般性介绍。

\section{数据自适应参数的 CV-TMLE}

令 $\mathcal D$ 为参数所有可能定义所组成集合的指标集。例如，在我们的情形中，它就是所有 $(a_L,a_H)$ 取值的集合。另外，假设对每个 $d\in\mathcal D$，都有一个统计目标参数 $\Psi_d:\mathcal M\to\mathbb R$。例如，若 $d$ 表示某个处理规则，可以定义 $\Psi_d(P)=E_PY_d$。令 $\hat d:\mathcal M_{NP}\to\mathcal D$ 为将经验分布映射到所需指标 $d_0$ 的一个估计的算法。在我们的例子中，$\hat d(P_n)=(a_L(P_n),a_H(P_n))$，而在 $P_{n,v^c}$ 上学得的相应目标参数 (9.8) 可写为
\[
\begin{aligned}
\Psi_{\hat d(P_{n,v^c})}(P)
={}&E_P\bigl\{E_P(Y\mid A=a_H(P_{n,v^c}),W)\\
&\hspace{2.6em}-E_P(Y\mid A=a_L(P_{n,v^c}),W)\bigr\}.
\end{aligned}
\]
本章着眼于介绍一种方法，为以下由 $\hat d$ 指标化的数据自适应目标参数（如式 (9.6)）提供估计量和统计推断：
\[
\psi_{0,n}=\operatorname{Ave}\{\Psi_{\hat d(P_{n,v^c})}(P_0)\}
=\frac1V\sum_{v=1}^V\Psi_{\hat d(P_{n,v^c})}(P_0).
\]
下面我们对前述 TMLE 估计量的平均作出修改，所得方法称为 CV-TMLE，记为 $\psi_n^*$。已有结果（van der Laan 和 Luedtke，2014）表明，这一 CV-TMLE $\psi_n^*$ 能提供稳健统计推断，而不依赖于会限制 $\hat d(P_n)$ 及 $\Psi_{\hat d(P_{n,v^c})}(P_0)$ 的相应估计量自适应程度的经验过程条件，即 Donsker 类条件。

对每个目标参数 $\Psi_d$，令 $D_d^*(P_0)$ 为其在 $P_0$ 处的有效影响曲线。假设 $\Psi_d(P_0)=\Psi_d(Q_0^d)$ 仅通过参数 $Q_0^d$ 依赖于 $P_0$，并假设 $D_d^*(P_0)=D_d^*(Q_0^d,g_0^d)$ 通过 $Q_0^d$ 和干扰参数 $g_0^d$ 依赖于 $P_0$。这些干扰参数以 $d$ 为指标，是因为目标参数的选择 $d$ 会影响这些参数的定义。

$\Psi_d:\mathcal M\to\mathbb R$ 的路径导数的典范梯度 $D_d^*(P)$ 蕴含一个二阶 Taylor 展开，其二阶项为 $R_d(Q^d,Q_0^d,g^d,g_0^d)$：
\[
\Psi_d(Q^d)-\Psi_d(Q_0^d)
=(P-P_0)D_d^*(Q^d,g^d)+R_d(Q^d,Q_0^d,g^d,g_0^d).
\]
令 $\hat Q^d:\mathcal M_{NP}\to\mathcal Q^d$ 和 $\hat g^d:\mathcal M_{NP}\to\mathcal G^d$ 分别为 $Q_0^d$ 和 $g_0^d$ 的初始估计量；$L^d(Q^d)$ 是 $Q_0^d$ 的一个有效损失函数，满足 $Q_0^d=\argmin_{Q^d}P_0L^d(Q^d)$；$\{Q^d(\epsilon):\epsilon\}$ 是在 $\epsilon=0$ 时经过 $Q^d$ 的子模型，具有一维或多维参数 $\epsilon$，使得广义得分的线性包包含 $(Q^d,g^d)$ 处的有效影响曲线：
\[
D_d^*(Q^d,g^d)\in
\left\langle\left.\frac{d}{d\epsilon}L^d(Q_\epsilon^d)\right|_{\epsilon=0}\right\rangle,
\]
其中，$\langle f\rangle=\{\sum_j\beta_j f_j:\beta\}$ 表示由 $f$ 的各个分量张成的线性空间。对经验分布为 $P_n$ 的样本，令 $\{\hat Q_\epsilon^d(P_n):\epsilon\}$ 为使用 $\hat g^d(P_n)$、在 $\epsilon=0$ 时经过估计量 $\hat Q^d(P_n)$ 的上述子模型。考虑 TMLE 只需一步的情形；采用通用最不利子模型 $\{\hat Q_\epsilon^d:\epsilon\}$，便可在形式上确保这一点。我们定义
\[
\epsilon_{n,v}=\argmin_{\epsilon}P_{n,v}L^d(Q^d_{n,v^c,\epsilon}),
\]
其中，子模型 $\{Q^d_{n,v^c,\epsilon}:\epsilon\}$ 在 $\epsilon=0$ 时经过 $Q^d_{n,v^c}=\hat Q^d(P_{n,v^c})$，并使用 $g^d_{n,v^c}=\hat g^d(P_{n,v^c})$ 作为 $g_0^d$ 的估计量。由此定义了基于 $P_{n,v}$、对 $Q^d_{n,v^c}$ 的第一步 TMLE 更新。如果考虑 TMLE 一步即收敛的情形，那么这一第一步 TMLE 就已经求解了其有效影响曲线得分方程：采用通用最不利子模型总可以实现一步收敛，而我们的例子也属于这种情形。
\[
P_{n,v}D_d^*(Q^d_{n,v^c,\epsilon_{n,v}},g^d_{n,v^c})=0.
\]
由于它对每个 $v$ 都成立，因此还有
\begin{equation}
\frac1V\sum_{v=1}^V P_{n,v}D_d^*(Q^d_{n,v^c,\epsilon_{n,v}},g^d_{n,v^c})=0.
\tag{9.16}
\end{equation}
关键方程 (9.16) 就是目标参数所需的有效得分方程。在 TMLE 只需一步的这一情形中，也可以通过一个对所有 $v$ 共用的 MLE 来满足该关键得分方程：
\[
\epsilon_n=\argmin_{\epsilon}\frac1V\sum_{v=1}^V P_{n,v}L^d(Q^d_{n,v^c,\epsilon}).
\]
一般而言，如果 TMLE 由多步 TMLE 算法定义，则针对每个 $v$，分别将该多步 TMLE 算法应用于 $P_{n,v}$，以确定 $Q^d_{n,v^c}$ 的 TMLE 更新 $Q^{d*}_{n,v^c}$；同样，也可以在这一多步 TMLE 算法的每一步汇集各个 $v$ 的信息。可以将 TMLE 更新 $Q^{d*}_{n,v^c}$ 看作基于 $P_{n,v}$ 的 TMLE 更新，只是其初始估计量 $Q^d_{n,v^c}$ 基于外部样本 $P_{n,v^c}$。

为简化记号，写 $Q_{n,v^c}=\hat Q^{\hat d(P_{n,v^c})}(P_{n,v^c})$，并类似地定义 $g_{n,v^c}=\hat g^{\hat d(P_{n,v^c})}(P_{n,v^c})$。下文假设采用一步 TMLE，但将其推广到迭代 TMLE 是直接的。

对 $\epsilon_n$ 及相应更新 $Q_{n,v^c,\epsilon_n}$ 的关键假设，是它能够求解有效影响曲线的交叉验证经验均值方程：
\begin{equation}
\frac1V\sum_{v=1}^V P_{n,v}
D^*_{\hat d(P_{n,v^c})}(Q_{n,v^c,\epsilon_n},g_{n,v^c})
=o_P(1/\sqrt n).
\tag{9.17}
\end{equation}
如果采用完整的 TMLE 更新，那么如上所示，将 $o_P(1/\sqrt n)$ 替换为 0 后，该方程仍成立；如果采用的 $Q_0^d,g_0^d$ 的初始估计量收敛速度快于 $n^{-1/4}$，则可证明，在很一般的情况下，第一步 TMLE 仍满足式 (9.17)。

所提的 $\psi_{0,n}$ 估计量为
\[
\psi_n^*\equiv\frac1V\sum_{v=1}^V
\Psi_{\hat d(P_{n,v^c})}(Q_{n,v^c,\epsilon_n}).
\]
在现有文献中，我们将这一估计量称为 CV-TMLE。相对于原始 CV-TMLE，唯一的变化是：在每个训练样本上，我们将目标改为由该训练样本上的 $\hat d(P_{n,v^c})$ 所隐含的、特定于该训练样本的目标参数；而在原始 CV-TMLE 表述中，目标始终是 $\Psi_d(P_0)$。经过这一微小变化，（同一个）CV-TMLE 如今用于靶向特定于训练样本的目标参数在 $V$ 个训练样本间的平均。这一 CV-TMLE 的一般渐近定理见第 A.1 节。

假设 $g_0^d$ 已知，并使用其已知值，从而 $\hat g^d(P_n)=g_0^d$。考虑下列估计量
\begin{equation}
\sigma_n^2=\underset{v}{\operatorname{Ave}}\left[
P_{n,v}\left\{D^*_{\hat d(P_{n,v^c})}(Q_{n,v^c,\epsilon_n},g_{n,v^c})\right\}^2\right],
\tag{9.18}
\end{equation}
它估计的是 CV-TMLE $\psi_n^*$ 的渐近方差 $\sigma_0^2=P_0\{D^*_{d_0}(Q,g_0)\}^2$。换言之，它是在各个 $v$ 之间，对特定于 $v$ 的影响曲线样本方差取平均；该影响曲线对应于以估计样本 $P_{n,v}$ 为基础、用于估计 $\Psi_{\hat d(P_{n,v^c})}(P_0)$ 的特定于 $v$ 的 CV-TMLE $\Psi_{\hat d(P_{n,v^c})}(Q_{n,v^c,\epsilon_n})$。由于计算特定于 $v$ 的有效影响曲线的代入估计所需的估计值 $(Q_{n,v^c,\epsilon_n},g_{n,v^c})$，本来就已用于估计量 $\psi_n^*$，所以这一 $D^*_{d_0}$ 方差估计不增加计算负担。最后，根据第 A.1 节的结果，$\psi_{0,n}$ 的一个渐近 0.95 置信区间为 $\psi_n^*\pm1.96\sigma_n/\sqrt n$。如果 $g_0^d$ 未知，而 $\hat Q^d$ 和 $\hat g^d$ 分别为 $Q_0^d$ 和 $g_0^d$ 的一致估计量，该置信区间也渐近有效。当 $g_0$ 未知，且 $\hat g(P_n)$ 是依据某个正确设定模型得到的 $g_0$ 的 MLE 时，也可以采用同一方差估计量和置信区间。此时，定理告诉我们，若 $\hat Q^d$ 不一致，该置信区间在渐近意义下是保守的。无论哪种情况，只要能够依赖 $\hat g^d$ 的一致性，我们一般都推荐采用这一置信区间。
